Reverberation mapping (or echo mapping) is a technique used to estimate the
mass of a black hole at the center of an Active Galactic Nucleus by observing
the variations in light emitted from the system. In this setup, a repeating
dot flash represents a sudden flare from the accretion disk, and a stationary
concentric ring appearing after a while represents the light echoing through
the gas region around the central black hole.

\begin{description}
\item[(a)] (5 points) Target Acquisition: Point your telescope toward the top
  floor of the Phenikaa Tower (the highest building on the Phenikaa campus),
  where the central dot and ring flashes will appear. Call your examiner for
  confirmation and marking. (If you cannot find the target, you may ask your
  examiner to point the telescope for you for 0 points in this part.)
\item[(b)] (10 points) Measure the time delay $(\tau)$ between the dot flash
  and the subsequent ring at least 3 times (up to 5 times). Record all
  individual readings, and calculate the mean time delay and its uncertainty.
\item[(c)] (5 points) Assume the gas orbiting the black hole with velocity of
  $3000\,\mathrm{km\,s^{-1}}$ and one measured second corresponding to one
  day. Calculate the black hole's mass $M$ (in $\mathrm{M}_\odot$).
\end{description}
